\section{Discussion and Limitations}
\label{sec:discussion}

\paragraph{The uniform-weight baseline is hard to beat here.} Uniform WRRR has a closed form and
does not concentrate on the task that looks worst on the calibration sample. On fresh draws it was
not worse than the empirical minimax solution on average, and its calibration optimism was smaller
(Table~\ref{tab:optimism}). We therefore keep it as the primary baseline for any further work. A
worst-task objective would have to beat it at the same rank and the same calibration access
before it is worth its cost. In stage~2 a minimax fit took a median of \STwoMinimaxCPUMedian~s
of CPU time (at most \STwoMinimaxCPUMax~s, including the bound audit), against a median of
\STwoUniformCPUMedianMs~ms for the closed form.

\paragraph{What we could not separate.}
\begin{itemize}
\item \emph{Relative normalisation.} Uniform WRRR and minimax both use relative errors. The
relative-error choice is itself a design decision, and we did not compare it with absolute errors
in this study. Its effect is not separated from the effects reported here.
\item \emph{Moment--normaliser interaction.} The Emp/Pop arm was worse than Emp/Emp. A plausible
reading is that errors in $\Shat_k$ and in $\dhat_k$, which are computed from the same draw, partly
cancel in the ratio. We did not test this.
\item \emph{Cause.} The pre-registered rule leaves the cause undetermined. The oracle arms tell us
what changes when an input to the objective is replaced by its population value. They do not
identify a mechanism, and population moments are not available to a deployable method.
\end{itemize}

\paragraph{Statistical limitations.} The design has \CfgNTeachers{} teachers and \CfgNResamples{}
calibration draws nested in each teacher. We report means, spreads and per-draw values
(Table~\ref{tab:perdraw}) and run no significance test. A non-significant or ambiguous difference
is not evidence that the two methods are equivalent. The original draw was seen before the
diagnostic was designed, so stage~2 is a development diagnostic, not a confirmatory experiment.

\paragraph{Scope of the setting.} The objective is the error of one linear layer's output. It is
not an end-to-end task loss, and gains in it need not transfer to accuracy. The synthetic
generator uses Gaussian inputs with a known second moment and a single configuration of dimension,
rank and sample size. We ran no sweep over these quantities. The bounds of
Section~\ref{sec:analysis} are deterministic and hold on a norm-bounded analysis set, whereas the
solver optimises over the unbounded rank-constrained set.

\paragraph{Real adapters.} We checked \AdmNAdapters{} public LoRA adapters (Flan-T5-base, GLUE CoLA, MRPC,
RTE and SST-2; \citealp{chung2022flan,tang2024fusionbench}) using metadata only. They were excluded
for three reasons:
\begin{itemize}
\item the adapter repositories carry no license information;
\item the base-model revision used for training is not recorded;
\item the training and checkpoint-selection split is not documented.
\end{itemize}
The publisher evaluates on the GLUE validation split, which is the split we had planned to use for
development and testing. The publisher also reports exact-match accuracy of generated label words
for all four tasks. The official GLUE metrics differ for two of them: Matthews correlation for
CoLA and accuracy/F1 for MRPC \citep{wang2018glue}. These scales should not be
mixed into a single percentage-point gain. Appendix~\ref{app:planned} lists the planned but
unexecuted real-adapter pilot.
